\part*{Parte A: Formulación Teórica del Problema}
\addcontentsline{toc}{section}{Parte A: Formulación Teórica del Problema}

\section{Introducción: El Mundo de los Fluidos en Movimiento}
Habiendo dominado problemas de difusión y mecánica de sólidos, nos adentramos en una de las disciplinas más desafiantes y valiosas de la simulación: la **Dinámica de Fluidos Computacional (CFD)**. Nuestro objetivo es modelar y predecir el movimiento de fluidos, un campo con aplicaciones en casi todas las ramas de la ingeniería y la ciencia.

Comenzaremos nuestro viaje con el caso más simple y fundamental: el **Flujo de Stokes**. Esta es la descripción del movimiento de fluidos donde las fuerzas viscosas (la ``fricción interna'' del fluido) dominan completamente sobre las fuerzas inerciales (la ``tendencia del fluido a seguir en movimiento''). Esto ocurre en flujos muy lentos o con fluidos extremadamente viscosos. El Flujo de Stokes es el ``Hola, Mundo'' de la CFD, ya que nos introduce al acoplamiento de los campos de velocidad y presión sin la complejidad añadida de la no linealidad.

\begin{infobox}
    Aunque es una simplificación, el Flujo de Stokes describe con gran precisión una enorme variedad de fenómenos reales:
    \begin{itemize}
        \item \textbf{Geofísica:} El movimiento del manto terrestre, que se comporta como un fluido extremadamente viscoso durante escalas de tiempo geológicas.
        \item \textbf{Microfluídica:} El flujo en dispositivos de ``laboratorio en un chip'' (lab-on-a-chip), donde las pequeñas dimensiones hacen que las fuerzas viscosas sean dominantes.
        \item \textbf{Ingeniería de Procesos:} El flujo de polímeros fundidos, vidrios o metales líquidos durante su procesamiento.
        \item \textbf{Bioingeniería:} El movimiento de microorganismos como bacterias o el flujo de sangre en capilares muy pequeños.
    \end{itemize}
\end{infobox}

\section{La Física de un Fluido Viscoso e Incompresible}
El movimiento de cualquier fluido newtoniano se describe por las célebres \textbf{ecuaciones de Navier-Stokes}. Para un fluido de densidad $\rho$ y viscosidad dinámica $\mu$ constantes, estas ecuaciones son:

\begin{align}
    \rho \left( \frac{\partial \mathbf{u}}{\partial t} + (\mathbf{u} \cdot \nabla) \mathbf{u} \right) &= -\nabla p + \mu \nabla^2 \mathbf{u} + \mathbf{f} && \text{(Conservación de Momento)} \\
    \nabla \cdot \mathbf{u} &= 0 && \text{(Conservación de Masa / Incompresibilidad)}
\end{align}

Aquí, $\mathbf{u}$ es el campo de velocidad, $p$ es el campo de presión y $\mathbf{f}$ son las fuerzas externas.

\subsection{La Simplificación de Stokes: El Régimen de Bajo Reynolds}
Las ecuaciones de Navier-Stokes son notoriamente difíciles de resolver debido al término no lineal de convección $(\mathbf{u} \cdot \nabla) \mathbf{u}$. Sin embargo, en muchos sistemas, este término es despreciable. El criterio para esta simplificación es el **número de Reynolds ($Re$)**:
$$ Re = \frac{\rho U L}{\mu} $$
Este número adimensional compara la magnitud de las fuerzas de inercia (numerador) con las fuerzas viscosas (denominador). Cuando $Re \ll 1$, el flujo es laminar y dominado por la viscosidad.

Bajo esta condición, podemos simplificar las ecuaciones. Asumiendo además un estado estacionario ($\frac{\partial \mathbf{u}}{\partial t} = 0$) y sin fuerzas externas ($\mathbf{f}=\mathbf{0}$), el término convectivo desaparece y las ecuaciones de Navier-Stokes se reducen a las \textbf{Ecuaciones de Stokes}:
\begin{align}
    -\mu \nabla^2 \mathbf{u} + \nabla p &= \mathbf{0} && \text{(Equilibrio de Momento)} \label{eq:stokes_mom} \\
    \nabla \cdot \mathbf{u} &= 0 && \text{(Incompresibilidad)} \label{eq:stokes_inc}
\end{align}
Este es un sistema de EDPs \textit{lineal} para las incógnitas $\mathbf{u}$ y $p$.

\section{La Forma Débil: Un Problema de Punto de Silla}
A diferencia de los módulos anteriores donde teníamos una única incógnita, aquí tenemos un sistema acoplado para la velocidad $\mathbf{u}$ y la presión $p$. Esto da lugar a lo que se conoce como un \textbf{problema de punto de silla}.

Para derivar la forma débil, multiplicamos cada ecuación por su correspondiente función de prueba e integramos:
\begin{enumerate}
    \item Multiplicamos la Ecuación \ref{eq:stokes_mom} por una función de prueba vectorial $\mathbf{v}$ e integramos sobre $\Omega$.
    \item Multiplicamos la Ecuación \ref{eq:stokes_inc} por una función de prueba escalar $q$ e integramos sobre $\Omega$.
\end{enumerate}
$$ \int_{\Omega} (-\mu \nabla^2 \mathbf{u} + \nabla p) \cdot \mathbf{v} \, dV = 0 \quad \text{y} \quad \int_{\Omega} (\nabla \cdot \mathbf{u}) q \, dV = 0 $$
Aplicando integración por partes al término de la viscosidad y al de la presión en la primera ecuación, obtenemos la forma débil final.

\begin{keybox}{WasaffCopper}
    \textbf{Sistema Variacional de Stokes:}
    
    Encontrar el par ($\mathbf{u}, p$) tal que para todo par de funciones de prueba ($\mathbf{v}, q$):
    \begin{align}
        \int_{\Omega} \mu \nabla\mathbf{u} : \nabla\mathbf{v} \, dV - \int_{\Omega} p (\nabla \cdot \mathbf{v}) \, dV &= \int_{\partial\Omega} (\mu \nabla\mathbf{u} - p\mathbf{I})\mathbf{n} \cdot \mathbf{v} \, dS \\
        \int_{\Omega} q (\nabla \cdot \mathbf{u}) \, dV &= 0
    \end{align}
\end{keybox}

\begin{notatecnica}
    \textbf{La Condición de Estabilidad LBB (inf-sup):}
    
    Un aspecto crucial en los problemas de Stokes es que no cualquier combinación de espacios de funciones para $\mathbf{u}$ y $p$ funciona. Si se eligen espacios "incompatibles" (por ejemplo, elementos lineales P1 para ambos), la solución de la presión puede presentar oscilaciones espurias y sin sentido físico.
    
    Para garantizar una solución estable y única, los espacios de funciones deben satisfacer la condición de estabilidad de Ladyzhenskaya-Babuška-Brezzi (LBB). En la práctica, esto significa que el espacio para la velocidad debe ser ``más rico'' que el espacio para la presión. La elección más común y robusta, conocida como el elemento **Taylor-Hood**, utiliza elementos cuadráticos (P2) para la velocidad y elementos lineales (P1) para la presión. Esta será nuestra elección en la implementación.
\end{notatecnica}
